% Independent draft sections for integration, not a standalone manuscript.
% Suggested citation keys: li2020, fanaskov2022, efendiev2013.
% Quantitative sources are documented in medium_review.md.

\section{Introduction}
Repeated scattering calculations require the wave field at several frequencies and incident directions. A reduced approximation can store spatial information in a fixed basis and replace repeated solution of the full finite element system by prediction of a small coefficient vector. This separates two questions: whether the basis represents the fields accurately, and whether the learned map predicts the coefficients needed for that accuracy. These questions need not have the same answer.

We study this distinction for a two-dimensional Helmholtz scattering problem with one fixed heterogeneous medium. A plane wave with wave number $k$ and direction angle $\theta$ produces the complex scattered field $u_s$. We fit separate proper orthogonal decomposition (POD) bases to the real and imaginary parts of training fields, then learn a coefficient map from $(k,\cos\theta,\sin\theta)$. The model is a POD adaptation of coefficient-space spectral neural operators. Since the material does not change in the main experiment, the data test a parametric solution family rather than an operator over varying material fields.

The main experiment uses one model, one wave-number normalization, and common output bases over $k\in[0.1,4]$. We compare its predictions with four prescribed generalized multiscale finite element method (GMsFEM) configurations on the same 424 held-out queries and a common reference finite element solution. We also retain two separate diagnostics: a low-frequency comparison of POD, Fourier, and Chebyshev representations on another fixed-medium dataset, and a frozen low-band model tested outside its training frequency range. Together, these records allow us to distinguish accurate interpolation within a sampled family from frequency extrapolation.

\section{Related work}
Fourier neural operators parameterize learned integral operators in Fourier space~\cite{li2020}. Spectral neural operators instead formulate learned maps directly between series coefficients, with Fourier and Chebyshev representations among the constructions studied by Fanaskov and Oseledets~\cite{fanaskov2022}. Our model uses data-dependent POD modes on an unstructured finite element mesh and a dense learned mixing of coefficient positions. It is therefore an adaptation of coefficient-space learning; the modes are neither Fourier frequencies nor eigenfunctions of the Helmholtz operator. The combination of a reduced basis with a learned coefficient map belongs to established nonintrusive reduced modeling, and is not itself a new principle.

GMsFEM constructs local approximation spaces and couples them through a global coarse problem~\cite{efendiev2013}. Here it provides a physics-based comparison with the learned global representation. The prescribed local spaces use 16 modes per coarse vertex, whereas the POD ranks count global modes for each real component. These counts describe different approximation spaces and do not define an equal computational budget. Our comparison concerns the recorded settings and reference discretization; it does not estimate the best accuracy attainable by either family.

\section{Interpretation}
On the common test set, the shared model has mean relative complex $L^2$ errors of $2.3656\%$ for the total field and $3.3271\%$ for the scattered field. The corresponding means for the best of the four recorded GMsFEM configurations, a $20\times20$ coarse grid with Type~2 snapshots and 16 local modes, are $3.2439\%$ and $4.2925\%$. Each relative error uses the norm of its own reference field as denominator. Thus, the two reported percentages answer different questions even though adding the same analytic incident wave leaves the complex error field unchanged.

The lower mean does not imply lower error for every query. The shared model has a total-field 95th percentile of $7.3718\%$ and a maximum of $27.5856\%$. Its total-field median is $1.3316\%$, while the selected GMsFEM configuration has a median of $1.1848\%$. The interval summaries also change ordering: GMsFEM has smaller mean errors in the two reported intervals up to $k=2$, and the shared model has smaller means in the two higher intervals. The comparison depends on the sampled frequency and direction distribution.

The shared model's fixed affine decoder admits a mean test scattered-field error of $0.1435\%$ under the best consistent-mass projection into its span. This floor is much smaller than the observed prediction error. For this run, decoder capacity alone therefore does not explain the remaining error. This comparison does not identify its cause: optimization, training coverage, normalization, and the coefficient map have not been separated by controlled ablations. The orthogonal error identity applies to squared errors for individual fields; subtracting mean relative errors does not recover a coefficient-error term.

The separate static-low study has a different balance. There, the POD model's mean test scattered-field error is $0.9475\%$, close to its $0.9188\%$ consistent-decoder floor. The recorded FNO and Chebyshev-SNO means are $1.6599\%$ and $7.9364\%$. For the Chebyshev model, a validation-only projection into the actual $8\times8$ decoder span has mean error $7.6143\%$. This diagnostic shows that the chosen polynomial span already limits reconstruction accuracy. Since the models use different representations, capacities, and devices, these records do not establish a general ranking of the architecture families.

The frozen model trained on $k\in[0.1,2]$ fails on the held-out queries with $k>2$: its mean total-field error is $131.8231\%$, and its mean scattered-field error is $177.7174\%$. This model has its own low-band basis and normalization. The result concerns that frozen surrogate and does not isolate basis error from coefficient-map error. The shared model is trained across $[0.1,4]$, so its high-frequency results are interpolation within the declared range rather than evidence of successful extrapolation.

\section{Limitations}
Both completed datasets contain one fixed material realization. Splitting by frequency prevents the same physical frequency from entering both training and test, but it does not test new material configurations, new meshes, or a different boundary model. The 424 main test fields contain 52 frequency groups, so fields at different directions are not 424 independent frequency samples. We report one training seed and provide no estimate of variability across repeated training runs. The selected checkpoint reaches the prescribed epoch limit; convergence to the best attainable predictor is not established.

The main reference is a P1 finite element solution on the saved $h=0.02$ mesh. GMsFEM starts from a different $h=0.04$ fine mesh, whose mean total-field discrepancy from the reference is $0.5322\%$. The GMsFEM comparison therefore contains both reduction error and the difference between these discretizations. All methods are sampled at the same reference quadrature points, with the incident wave evaluated analytically. This protocol does not remove finite element discretization error, the error of the absorbing boundary condition, or possible quadrature error when fields from different meshes are compared.

POD fitting uses a diagonal lumped mass weighting, whereas training measures the consistent finite element norm through the decoder Gram matrices and the irreducible residual. The resulting loss is exact for the fixed decoder and its specified discrete targets, up to the stated numerical precision; this is not a claim that the fitted POD is optimal in the consistent mass norm. The comparison also gives neither an equal-budget benchmark nor an end-to-end speedup: supervised label generation, frequency-dependent multiscale setup, and heterogeneous hardware must be accounted for before such a claim can be tested.
